% ============================================================
% style/boxes.tex
%
% Универсальные декоративные рамки для документа.
% ============================================================

% ------------------------------------------------------------
% 1. Основные параметры рамок
% ------------------------------------------------------------

\newcommand{\rboxlinewidth}{0.75pt}
\newcommand{\rboxlinecolor}{rframe}


% ------------------------------------------------------------
% 2. Рисование декоративной рамки
% ------------------------------------------------------------
%
% Типы:
%   {[}    рамка без правой стороны
%   {[]}   полный прямоугольник
%   {^}    рамка без нижней стороны
%   {v}    рамка без верхней стороны
%
%   {<u}   левый верхний угловой контур
%   {<d}   левый нижний угловой контур
%   {<}    то же самое, что {<d}
%
%   {>u}   правый верхний угловой контур
%   {>d}   правый нижний угловой контур
%   {>}    то же самое, что {>d}

\newcommand{\rboxdraw}[1]{%
    \ifstrequal{#1}{[}{%
        \draw[\rboxlinecolor,line width=\rboxlinewidth]
            (frame.north east)
            -- (frame.north west)
            -- (frame.south west)
            -- (frame.south east);%
    }{}%
    \ifstrequal{#1}{[]}{%
        \draw[\rboxlinecolor,line width=\rboxlinewidth]
            (frame.north west)
            rectangle
            (frame.south east);%
    }{}%
    \ifstrequal{#1}{^}{%
        \draw[\rboxlinecolor,line width=\rboxlinewidth]
            (frame.south west)
            -- (frame.north west)
            -- (frame.north east)
            -- (frame.south east);%
    }{}%
    \ifstrequal{#1}{v}{%
        \draw[\rboxlinecolor,line width=\rboxlinewidth]
            (frame.north west)
            -- (frame.south west)
            -- (frame.south east)
            -- (frame.north east);%
    }{}%

    % Левый верхний угол: левая сторона + верхняя сторона
    \ifstrequal{#1}{<u}{%
        \draw[\rboxlinecolor,line width=\rboxlinewidth]
            (frame.south west)
            -- (frame.north west)
            -- (frame.north east);%
    }{}%

    % Левый нижний угол: левая сторона + нижняя сторона
    \ifstrequal{#1}{<d}{%
        \draw[\rboxlinecolor,line width=\rboxlinewidth]
            (frame.north west)
            -- (frame.south west)
            -- (frame.south east);%
    }{}%

    % Старый вариант: то же самое, что <d
    \ifstrequal{#1}{<}{%
        \draw[\rboxlinecolor,line width=\rboxlinewidth]
            (frame.north west)
            -- (frame.south west)
            -- (frame.south east);%
    }{}%

    % Правый верхний угол: верхняя сторона + правая сторона
    \ifstrequal{#1}{>u}{%
        \draw[\rboxlinecolor,line width=\rboxlinewidth]
            (frame.north west)
            -- (frame.north east)
            -- (frame.south east);%
    }{}%

    % Правый нижний угол: нижняя сторона + правая сторона
    \ifstrequal{#1}{>d}{%
        \draw[\rboxlinecolor,line width=\rboxlinewidth]
            (frame.south west)
            -- (frame.south east)
            -- (frame.north east);%
    }{}%

    % Старый вариант: то же самое, что >d
    \ifstrequal{#1}{>}{%
        \draw[\rboxlinecolor,line width=\rboxlinewidth]
            (frame.south west)
            -- (frame.south east)
            -- (frame.north east);%
    }{}%
}


% ------------------------------------------------------------
% 3. Универсальная среда rbox
% ------------------------------------------------------------
%
% Использование:
%
%   \begin{rbox}{[]}
%       Текст внутри блока.
%   \end{rbox}
%
% С дополнительными настройками:
%
%   \begin{rbox}[left=24pt,right=24pt]{^}
%       Текст внутри блока.
%   \end{rbox}

\NewTColorBox{rbox}{O{} m}{%
    enhanced,%
    breakable,%
    width=\linewidth,%
    colback=rbg,%
    colframe=rbg,%
    coltext=rtext,%
    boxrule=0pt,%
    frame hidden,%
    sharp corners,%
    left=10pt,%
    right=10pt,%
    top=10pt,%
    bottom=10pt,%
    before skip=1.1em,%
    after skip=1.1em,%
    overlay={%
        \rboxdraw{#2}%
    },%
    #1%
}